% Şekil 3.1: genel sistem mimarisi. Her blokta işlevini gösteren küçük bir simge vardır.
% Düz oklar güç yolu, kesikli oklar ölçüm ve denetim yoludur. Yalnız TikZ kullanır.
\begin{tikzpicture}[x=1mm,y=1mm,
    kutu/.style={draw, rounded corners=2pt, minimum width=28mm, minimum height=26mm},
    guc/.style={kutu, fill=blue!5},
    den/.style={kutu, fill=orange!10},
    olc/.style={kutu, fill=gray!10},
    ad/.style={font=\footnotesize, align=center},
    simge/.style={line width=0.6pt, line cap=round, line join=round},
    flow/.style={-{Latex[length=2mm]}, thick},
    sig/.style={-{Latex[length=2mm]}, thick, dashed},
    nokta/.style={circle, fill=black, inner sep=0.9pt}
  ]
  % ---------------------------------------------------------------- bloklar
  \node[guc] (bat)  at (0,0)    {};
  \node[guc] (dc)   at (38,0)   {};
  \node[guc] (inv)  at (76,0)   {};
  \node[guc] (lcl)  at (114,0)  {};
  \node[guc] (grid) at (152,0)  {};
  \node[den] (pwm)  at (76,-38) {};
  \node[olc] (meas) at (114,-38){};
  \node[den] (control) at (76,-72) {};
  \node[den] (sin3) at (38,-72) {};

  \node[ad] at (0,-8.5)    {Batarya /\\[1.2mm]DA kaynak};
  \node[ad] at (38,-8.5)   {Bölünmüş\\[1.2mm]DA bara};
  \node[ad] at (76,-8.5)   {T-tipi\\[1.2mm]güç katı};
  \node[ad] at (114,-8.5)  {LCL\\[1.2mm]filtre};
  \node[ad] at (152,-8.5)  {Şebeke /\\[1.2mm]ortak bağlantı};
  \node[ad] at (76,-46.5)  {PWM ve kapı\\[1.2mm]sinyalleri};
  \node[ad] at (114,-46.5) {Gerilim ve akım\\[1.2mm]ölçümleri};
  \node[ad] at (76,-80.5)  {Akım denetimi ve\\[1.2mm]sınırlandırıcılar};
  \node[ad] at (38,-80.5)  {{Sin\textsuperscript{3}}\\[1.2mm]kompanzasyonu};

  % ---------------------------------------------------------------- simgeler
  % Batarya: iki hücre
  \begin{scope}[shift={(0,4)}, simge]
    \draw (0,7) -- (0,2.4);
    \draw (-5,2.4) -- (5,2.4);
    \draw[line width=1.4pt] (-2.6,0.9) -- (2.6,0.9);
    \draw (-5,-0.8) -- (5,-0.8);
    \draw[line width=1.4pt] (-2.6,-2.3) -- (2.6,-2.3);
    \draw (0,-2.3) -- (0,-6.5);
    \node[font=\scriptsize] at (6.8,3.6) {$+$};
    \node[font=\scriptsize] at (6.8,-3.4) {$-$};
  \end{scope}

  % Bölünmüş DA bara: seri iki kapasitör ve orta nokta
  \begin{scope}[shift={(38,4)}, simge]
    \draw (0,7) -- (0,3.6);
    \draw (-4,3.6) -- (4,3.6);  \draw (-4,2.2) -- (4,2.2);
    \draw (0,2.2) -- (0,-2.2);
    \draw (-4,-2.2) -- (4,-2.2); \draw (-4,-3.6) -- (4,-3.6);
    \draw (0,-3.6) -- (0,-6.5);
    \draw (0,0) -- (7,0);
    \node[nokta] at (0,0) {};
    \node[font=\scriptsize, above] at (7,0) {$N$};
    \node[font=\scriptsize, left] at (-4,2.9) {$C_A$};
    \node[font=\scriptsize, left] at (-4,-2.9) {$C_B$};
  \end{scope}

  % T-tipi faz kolu: dış kolda iki anahtar, nötr kolda çift yönlü anahtar
  \begin{scope}[shift={(76,4)}, simge]
    % Şekil 3.5'teki anahtar kutusu gösterimiyle aynı
    \draw (4,7.5) -- (4,-7.5);
    \draw (-11,0) -- (10,0);
    \node[nokta] at (4,0) {};
    \foreach \x/\y/\q in {4/4/1, 4/-4/2, 0/0/3, -6/0/4}
      \node[draw, fill=white, minimum width=5mm, minimum height=3.6mm, inner sep=0pt,
            font=\scriptsize] at (\x,\y) {$Q_{\q}$};
    \node[font=\scriptsize, above] at (-10,0) {$N$};
  \end{scope}

  % LCL filtre: iki endüktans ve kapasitör
  \begin{scope}[shift={(114,4)}, simge]
    \draw (-11,2) -- (-9.5,2) arc(180:0:1) arc(180:0:1) arc(180:0:1) -- (3.5,2)
          arc(180:0:1) arc(180:0:1) arc(180:0:1) -- (11,2);
    \node[nokta] at (0,2) {};
    \draw (0,2) -- (0,-1.2);
    \draw (-3,-1.2) -- (3,-1.2); \draw (-3,-2.6) -- (3,-2.6);
    \draw (0,-2.6) -- (0,-4.6);
    \draw (-2.4,-4.6) -- (2.4,-4.6); \draw (-1.5,-5.5) -- (1.5,-5.5); \draw (-0.6,-6.4) -- (0.6,-6.4);
    \node[font=\scriptsize, above] at (-6.5,3) {$L_1$};
    \node[font=\scriptsize, above] at (6.5,3) {$L_2$};
    \node[font=\scriptsize, right] at (3,-1.9) {$C_f$};
  \end{scope}

  % Şebeke: sinüs kaynağı
  \begin{scope}[shift={(152,4)}, simge]
    \draw (-11,0) -- (-5.5,0);
    \draw (0,0) circle (5.5);
    \draw[domain=-3.6:3.6, samples=40, smooth] plot (\x, {2.2*sin(\x*50)});
  \end{scope}

  % PWM: taşıyıcı, referans ve kapı darbeleri
  \begin{scope}[shift={(76,-34)}, simge]
    \draw[line width=0.45pt] (-10,1.4) -- (-7.5,6.4) -- (-5,1.4) -- (-2.5,6.4) -- (0,1.4)
          -- (2.5,6.4) -- (5,1.4) -- (7.5,6.4) -- (10,1.4);
    \draw[red!70!black, line width=0.6pt] (-10.5,3.9) -- (10.5,3.9);
    \draw (-10,-1.2) -- (-8.75,-1.2) -- (-8.75,-5) -- (-6.25,-5) -- (-6.25,-1.2) -- (-3.75,-1.2)
          -- (-3.75,-5) -- (-1.25,-5) -- (-1.25,-1.2) -- (1.25,-1.2) -- (1.25,-5) -- (3.75,-5)
          -- (3.75,-1.2) -- (6.25,-1.2) -- (6.25,-5) -- (8.75,-5) -- (8.75,-1.2) -- (10,-1.2);
  \end{scope}

  % Ölçüm: gerilim ve akım ölçerleri
  \begin{scope}[shift={(114,-34)}, simge]
    \draw (-5,0.5) circle (3.6); \node[font=\footnotesize] at (-5,0.5) {V};
    \draw (5,0.5) circle (3.6);  \node[font=\footnotesize] at (5,0.5) {A};
  \end{scope}

  % Akım denetimi: toplama noktası, PI ve sınırlayıcı
  \begin{scope}[shift={(76,-68)}, simge]
    \draw (-8.5,0.5) circle (2);
    \node[font=\tiny] at (-8.5,0.5) {$\Sigma$};
    \draw[-{Latex[length=1.2mm]}] (-6.5,0.5) -- (-3.6,0.5);
    \draw (-3.6,-2) rectangle (2.4,3);
    \node[font=\scriptsize] at (-0.6,0.5) {PI};
    \draw[-{Latex[length=1.2mm]}] (2.4,0.5) -- (4.6,0.5);
    \draw (4.6,-2.2) rectangle (11,3.2);
    \draw[line width=0.7pt] (5.4,-1.2) -- (6.9,-1.2) -- (8.7,2.2) -- (10.2,2.2);
  \end{scope}

  % Sin3: sin^3 dalga biçimi, kesikli sinüs başvurusu
  \begin{scope}[shift={(38,-68)}, simge]
    \draw[gray, line width=0.4pt] (-10.5,0.5) -- (10.5,0.5);
    \draw[gray, dashed, line width=0.45pt, domain=-10:10, samples=60, smooth]
          plot (\x, {0.5+5*sin(\x*18)});
    \draw[blue!60!black, line width=0.8pt, domain=-10:10, samples=80, smooth]
          plot (\x, {0.5+5*sin(\x*18)*sin(\x*18)*sin(\x*18)});
  \end{scope}

  % ---------------------------------------------------------------- oklar
  \draw[flow] (bat) -- node[above, font=\scriptsize]{enerji} (dc);
  \draw[flow] (dc) -- (inv);
  \draw[flow] (inv) -- (lcl);
  \draw[flow] (lcl) -- (grid);
  \draw[sig] (lcl.south) -- (meas.north);
  \draw[sig] (grid.south) |- (meas.east);
  \draw[sig] (meas.south) |- (control.east);
  \draw[sig] (control) -- (pwm);
  \draw[sig] (pwm) -- (inv);
  \draw[sig] (sin3) -- (control);
\end{tikzpicture}%
